\documentclass[11pt]{article}
\usepackage[margin=1in]{geometry}
\usepackage{amsmath,amssymb,amsthm,booktabs,hyperref,xurl}
\newtheorem{theorem}{Theorem}
\newtheorem{lemma}{Lemma}
\newcommand{\F}{\mathbb F}\newcommand{\Hom}{\operatorname{Hom}}
\title{Acquiring Affine Hom Slices without a Common Original Center}
\author{Astra P3 constructive-representation research}
\date{27 October 2026}
\begin{document}\maketitle

\section{The obstruction removed}

The unsaturated module-evaluation counter can count an affine right slice
$T_0+W$ in a manufactured module-homomorphism space, but its preceding raw
frontend required all marked quadratic forms to share a single affine
bilinear center in the original independent inputs. These requirements
are different. An offset $T_0$ outside the actual right-input image $W$
cannot be removed by translating those original inputs.

The new frontend acquires such offsets from the actual linear terms of
the original quadratic registers. It does not receive a field,
decomposition, offset, module basis, or system of affine slice equations.
Every term of every original marked polynomial is retained. The code is
\path|affine_unsaturated_bilinear_counter.py|, with executed checks in
\path|affine_unsaturated_bilinear_audit.py|. The signed-fiber formula and
dual acquisition are the construction in
\path|UNSATURATED_MODULE_EVALUATION_THEOREM.tex|, which this note uses
without claiming a new formula or a general circuit algorithm.

\section{Acquired polar split and residual linear maps}

The input is an ordinary Boolean DAG with $n$ independent inputs and
$G$ explicit fan-in-two gates. The existing exact low-degree frontend
and complete-register reversible-suffix reversal acquire marked
quadratics $Q_1,\ldots,Q_m$ and one actual target. A successful joint
polar-adjoint certificate supplies complementary common isotropic
spaces and an invertible coordinate map $J$. Retain all residual
linear terms, so in these coordinates the complete map is
\begin{equation}
 Q(x_*+J(u,v))=B(u,v)+L(u)+R(v)+c.
 \label{eq:full}
\end{equation}
The polar construction depends only on the polar forms. It does not
need a common affine center. Its original default still requires that
center; the optional mode
\path|require_shared_center=False| exposes the residual maps in
\eqref{eq:full} and records whether a common original center existed.
The mode never deletes a linear monomial. Complete coefficient
identities are verified.

The prior nucleus, center, Frobenius, and sandwich-rank certificates
acquire a simple algebra $S\cong M_p(\F_{2^a})$ acting on $U,O$,
with $\dim_2U=apd$, $\dim_2O=apr$, and an injective map
\[
 \iota:V\longrightarrow H=\Hom_S(U,O),\qquad
 \iota(v)(u)=B(u,v).
\]
Its image is $W$, and its binary codimension is
$c_H=adr-\dim_2V$. No canonical finite-field or matrix-unit coordinates
are constructed.

\paragraph{Output dependencies are a real constraint.}
The output space $O$ is acquired from independent \emph{cross} forms.
A marked cross form can be dependent while its residual linear part is
independent. Such a mark cannot be treated as a constant dependency.
For every actual column of $L$ and $R$, the new code verifies every
acquired cross-output dependency. Only if all those identities hold
does it compress the linear maps to $O$. Independent constant offsets
remain in the full $m$-bit output constant. A linear-only violation is
declined explicitly; it is not silently projected away.

\section{Manufacturing the affine offset}

For each of the two acquired orientations of $U,V$, solve the actual
binary linear equations
\begin{equation}
 B(u_0,v)=R(v)\quad\hbox{for every basis vector }v\in V.
 \label{eq:right}
\end{equation}
If a solution exists it is unique, because the left radical of $B$ is
zero. Next express the actual matrix of $L:U\to O$ in the acquired basis
of $H$. This succeeds exactly when $L\in H$. Let $T_0=L$ as an element
of that space, and put $c'=c+L(u_0)$.

\begin{theorem}[Acquired affine evaluation]
Suppose the polar split, simple nucleus, residual output dependencies,
equations \eqref{eq:right}, and membership $L\in H$ have all been
verified. Then the complete original register map is
\[
 \boxed{Q(x_*+J(u'+u_0,v))
       =c'+\bigl(T_0+\iota(v)\bigr)(u').}
\]
This is a bijective reparameterization of the original $n$ input bits.
The affine offset $T_0$ is acquired by binary elimination even when
$T_0\notin W$. The construction introduces no independent variables.
\end{theorem}
\begin{proof}
Replace $u$ by $u'+u_0$ in \eqref{eq:full}. Bilinearity gives
$B(u',v)+B(u_0,v)+L(u')+L(u_0)+R(v)+c$.
The two terms $B(u_0,v)$ and $R(v)$ cancel in characteristic two by
\eqref{eq:right}. The remaining expression is exactly the displayed
formula. Translation on $U$ and the original map $J$ are bijections;
the right-radical test makes $V\to W$ injective and hence a bijection
onto its actual image. The program verifies the entire resulting ANF
of every marked output, including the new constant $c'$.
\end{proof}

If a common original affine center existed, the original routine would
find it and the residual maps would be zero. Conversely, in an admitted
orientation satisfying the theorem, translating the original right
input removes $T_0$ exactly when $T_0\in W$. Thus an acquired offset
outside $W$ records a substantive failure of the old shared-center
admission. The new implementation tests that membership rather than
assuming it.

\section{Exact counting and explicit costs}

The previous Hom-space construction manufactures $W^\perp$ and its
nondegenerate reverse trace pairing. For $y$ equal to the actual target
after constants and dependency checks, character orthogonality gives
\[
 \#\{(u',v):(T_0+\iota(v))(u')=y\}
 =2^{-c_H}\sum_{\lambda\in W^\perp}
   (-1)^{\lambda(T_0)}
   \sum_{T(u')=y}(-1)^{\lambda(T)}.
\]
The inner signed universal-evaluation count is the previously proved
polynomial formula in the acquired ranks of $y$, its reverse Hom
representative $D$, and $D(y)$. The affine frontend merely supplies the
now-manufactured phase $(-1)^{\lambda(T_0)}$ and the corrected original
output constant. The integer answer is verified to lie between zero
and $2^n$.

Let $N=n+m+G+2$ and let $L_{\rm in}$ be the ordinary explicit DAG
description length. On admitted inputs $N=O(L_{\rm in})$. The new
dependency tests, $u_0$ system, Hom-coordinate solve, and complete
quadratic reconstructions have polynomial binary size and fit inside
the existing conservative $O(N^{14})$ acquisition allowance. The
resulting deterministic bit bound is therefore still
\[
 O\bigl(N^{14}+2^{c_H}N^8\bigr),
\]
with polynomial space when characters are streamed. The count has at
most $n+1$ bits. This is polynomial for every fixed bound
$c_H\leq K\log_2 L_{\rm in}$. No such codimension bound is established
for arbitrary circuits, and this frontend is a verified sufficient
recognizer, not a complete recognizer of all equivalent computations.

The code checks both acquired orientations and chooses the smallest
codimension among those passing the full affine certificate. It does
not first choose an orientation from the cross forms alone and then
discard inconsistent linear terms. Its native guard rejects
$n>G+1$ before dense allocation, as required by zero common polar
radical and explicit fan-in-two ancestry. No original input assignment
is enumerated by the production counter.

\section{Executed evidence and remaining acquisition failure}

Five dense affine-hidden proper-slice fixtures over $\F_2$ and $\F_4$
passed every target comparison with literal original-input histograms:
$305664$ assignments and $4992$ targets in total. Their input counts
are $9,13,11,15,18$ and their binary codimensions are $1,2,1,1,2$.
Every fixture has a manufactured offset outside the actual right image;
the old shared-center frontend declines all five. Some slices have
binary codimension one inside an $\F_4$ Hom space and need not be
$\F_4$-linear.

Each fixture was also compiled into an ordinary gate DAG with 80
nonlinear reversible updates. The smaller programs were checked on
every original input and the larger ones on 256 original inputs.
Correct redundant linear-output dependencies and inconsistent
constant targets were tested. Adding a dependent cross output with a
new, nonzero linear term is declined in every case.

Two larger ordinary programs have 500 nonlinear reversible updates:
\begin{center}
\begin{tabular}{rrrrrr}
\toprule
$a$ & $(p,d,r)$ & $n$ & Gates & $c_H$ & Exact count\\
\midrule
1 & $(4,4,3)$ & 26 & 3327 & 2 & 23296\\
2 & $(3,2,3)$ & 22 & 3532 & 2 & 480\\
\bottomrule
\end{tabular}
\end{center}
The production consumer uses four characters in each. An independent
oracle enumerates 1024 actual right inputs and solves the original
field row equations for the left inputs; it uses neither the acquired
algebra nor the signed-fiber formula. Both counts agree, and both
programs pass 128 original gate evaluations. These executions establish
correctness, not a comparative speed claim.

The modified optional mode also passes the complete earlier
source-free-split audit with its default requirement retained.
Its earlier accepted family and exact centering identities are
unchanged.

The next obstruction is separate from the affine Hom formula. For
$X,Y\in\F_2^{2\times2}$, restrict the actual right input to
$Y_{11}=1$. The target $XY=0$ has 20 solutions among seven independent
input bits. The current joint-adjoint \emph{central} split nevertheless
declines this presentation because its ordinary center does not have
the required two Boolean atoms. The audit preserves this explicit
failure. It does not assert that every useful affine slice is acquired,
or that a different split-acquisition mechanism cannot recognize it.
Likewise no general sampler follows from the exact count; the
full-fiber rejection mechanism can be exponentially inefficient on
affine slices, as shown in the parent theorem.

The unrestricted polynomial-time Boolean counting goal remains
unproved. Classical module evaluation, the existing signed-fiber
formula, and the acquired affine interface are distinguished explicitly.
\end{document}
